% draft0: 結果 → 解釈 → 限界。実験節の数値を繰り返すだけにしない。
\section{考察}
\label{sec:discussion}

\subsection{RQへの回答と改善・失敗の解釈}
\begin{itemize}
    \item \textbf{書くこと：} RQ1は比較可能な完了手法の範囲，RQ2は旧ARの問題と現行構成の寄与，RQ3は順位の安定性と絶対校正の限界，という順に答える．
    \item \textbf{図案：} 改善・悪化の代表ケースで，観測系列／正常継続予測／残差／サービス順位を並べる．選定規則とcase IDを記し，良い例だけを選ばない（\hyperref[ev:main]{E3}・\hyperref[ev:history]{E5}）．
    \item \textbf{注意：} shift吸収と順位変化の関連は因果的証明ではない．CPUを含む障害種別差を隠さず，事後的な障害種別switchを提案しない．
\end{itemize}

\subsection{Bayes Factorの解釈とモデル仮定の限界}
\begin{itemize}
    \item \textbf{主張：} 正常内pseudo-faultでも大きな正のスコアが出るため，現方式は同一ケース内の相対順位付けとして解釈する（\hyperref[ev:robustness]{E7}）．
    \item \textbf{書くこと：} 正常時の非定常性，Gaussian・独立性仮定，多段誤差の残存相関，正常データの再利用，top-3集約の影響を切り分ける．
    \item \textbf{注意：} 正常holdout診断は「学習内分散の過小評価だけが原因」という説明を支持しない．原因が一意に判明したとは書かず，BF閾値・事後確率・case間比較へ拡大解釈しない．
\end{itemize}

\subsection{採用しなかった拡張と負の結果}
\begin{itemize}
    \item \textbf{書くこと：} BSRC-ARはAR生成regimeを直接比較する試みとして説明し，積分の収束と校正・RCA性能を分けて評価する．Adaptive Directは参考比較として位置付ける（\hyperref[ev:negative]{E8}）．
    \item \textbf{注意：} BSRC-ARを最終AMBERと混同しない．層化75ケースのgateと正式375ケースを分け，今回の具体化の負の結果をARとBayesの組合せ全般へ一般化しない．
    \item \textbf{TODO：} 本文は採用しない理由に絞り，詳細な仮説・積分・探索履歴は付録へ置く．この執筆ブランチでは再開発しない．
\end{itemize}

\subsection{妥当性への脅威と適用範囲}
\begin{itemize}
    \item \textbf{書くこと：} 同じRE1上での設計・次数探索，多重比較，ケース間の依存，少数システム・注入障害への依存を明記する．
    \item \textbf{注意：} 比較adapterの変更，欠測・空候補，時刻情報の利用，前処理の世代差，サービス当たりメトリクス数と集約の影響を公平性の観点で点検する．
    \item \textbf{範囲外：} 未知開始時刻の検出，実運用・外部benchmarkへの一般化，原因メトリクスの正解保証，因果方向の識別は今後の課題とする．
\end{itemize}
